% GENERATED by tools/gasp_partition_candidates.py from
%   data/gasp/partition_candidates.csv (the prose ledger) and
%   data/gasp/metadata_audit.csv (every number). Def A, the reach read and
%   the concordance rules are imported from tools/gasp_structural_baseline.py,
%   which emits the sibling audit table --- the 19-of-38 count appears in both
%   floats and in the chapter at subsec:attack-profiles, so it has one home.
% Do not hand-edit; regenerate. Requires booktabs (already in the preamble).
% Corpus at generation: 38 flows, 19/7/7/5, all high confidence.

\begin{table}[htbp]
\centering
\scriptsize
\setlength{\tabcolsep}{4pt}
\newcolumntype{R}[1]{>{\raggedright\arraybackslash}p{#1}}
\caption[Partition schemes considered]{Partition schemes considered for the objective-conditioned attack profiles, and what each produced on this corpus. Rows are ordered by how far the evidence a scheme reads sits from the analyst's own statement of what the operation did; the six above the rule were dismissed, and the row below it is the scheme the pipeline adopts. For the three schemes read from the incident's dependency graph, a technique is terminal when it has no outgoing edge to a different technique, and the disagreement counts are against the attested objective --- requiring an exact match, then in brackets allowing the two to share one objective tactic. A tick marks a scheme whose class membership could still be derived from artefacts this pipeline holds; an empty cell, one that cannot be built on this corpus at all. Attack Flow published export v3.1.1 (schema 2.0.0), 38 usable incidents, ATT\&CK Enterprise v19.1. Every count is computed from the classification audit at generation; the 47-campaign hand-label sample in the first row is the one figure reported from the project record.}
\label{tab:rejected-partitions}
\begin{tabular}{@{}R{0.150\textwidth}R{0.110\textwidth}cR{0.235\textwidth}R{0.310\textwidth}c@{}}
\toprule
Scheme & Slices on & Ref. & On this corpus & Why dismissed & Built? \\
\midrule
Motivation categories & Inferred motivation (STIX) &  & Ten categories, of which a hand-labelled sample of 47 external campaigns used three. The labels are written at campaign granularity and 0 of the 38 flows carry an ATT\&CK Campaign identifier to join them to. & Slices on what an analyst infers rather than on what the report records, and the labels do not join to this corpus at all. &  \\
\addlinespace
Group-witnessed objective & Actor attribution & \citep{mitre2026attackv19} & 18 of the 38 flows carry an ATT\&CK Group identifier and 0 carry a Campaign identifier; 20 are unattributed. & More than half the corpus is unclassifiable, and the classes that do form name actors rather than objectives. &  \\
\addlinespace
Terminal tactic in the dependency graph & Graph: terminal tactic &  & 7 exfiltration, 11 impact, 1 both, 19 neither. Disagrees with the attested objective on 19 of the 38 flows (15 where sharing one objective tactic counts as agreement). & The disagreement is systematic: 14 of the 19 are reports that stop before the objective is reached, so the graph is silent where the analyst attests one. It encodes documentation truncation as attacker behaviour. & \checkmark \\
\addlinespace
Objective tactic reached anywhere & Graph: objective reach &  & 10 exfiltration, 10 impact, 3 both, 15 neither, counting any occurrence rather than only the terminal one. Disagrees on 14 of the 38 (11 under the weaker rule). & More forgiving than the terminal read and still structural: 15 flows reach no objective tactic and stay unclassifiable. & \checkmark \\
\addlinespace
Terminal technique & Graph: terminal technique &  & 36 distinct terminal-technique sets across the 38 flows, 34 of them held by one flow alone; 2 flows are cyclic and have no terminal technique. & Fragments to roughly one class per incident. A class of one is a single account, not a behavioural envelope. Kept as a fragmentation control. &  \\
\addlinespace
Stated objective, three classes with multi-membership & Stated objective & \citep{alshamrani2019} & 26 exfiltration, 14 impact and 5 no-objective memberships --- 45 over 38 flows, because the 7 flows pursuing both objectives are counted twice. & The profiles stop being alternatives: a flow in two of them is weighed twice in the comparison, and the compound objective has no name of its own. & \checkmark \\
\midrule
\textbf{Stated objective, four disjoint classes} & Stated objective & \citep{alshamrani2019} & 19 exfiltration, 7 impact, 7 double extortion, 5 no realised objective. Every flow sits in exactly one class, all 38 at high audit confidence. & Adopted: the lightest scheme that is disjoint, sources every membership from an analyst's statement, and names the compound objective. &  \\
\bottomrule
\end{tabular}
\end{table}
